\chapter{Alessio Figalli --- Fields Medal 2018}\label{ch2018alessiofigalli:alessio-figalli-fields-medal-2018-a-comprehensive-report-on-his-work-and-impact}

\begin{center}\rule{0.5\linewidth}{0.5pt}\end{center}

\textbf{Prepared:} July 2026 \textbf{Subject:} Alessio Figalli (b. April
2, 1984, Rome, Italy) \textbf{Position:} Chaired Professor of
Mathematics, ETH Zurich; Director, Forschungsinstitut fur Mathematik
(FIM) \textbf{Doctoral Advisors:} Luigi Ambrosio (Scuola Normale
Superiore di Pisa), Cedric Villani (ENS Lyon) \textbf{Fields Medal
Citation:} \emph{``For contributions to the theory of optimal transport,
and its applications to partial differential equations, metric geometry,
and probability.''}

\begin{center}\rule{0.5\linewidth}{0.5pt}\end{center}

\section{Table of Contents}\label{ch2018alessiofigalli:table-of-contents}

\begin{enumerate}
\def\labelenumi{\arabic{enumi}.}
\tightlist
\item
 Introduction
\item
 Biographical Background
\item
 Optimal Transport Theory: An Overview
\item
 Regularity Theory for Optimal Transport Maps
\item
 The Monge--Ampere Equation
\item
 Partial Regularity for Monge--Ampere Type Equations
\item
 Quantitative Stability of Geometric and Functional Inequalities
\item
 The Anisotropic Isoperimetric Inequality
\item
 Stability of the Gagliardo--Nirenberg and Log-Sobolev Inequalities
\item
 Applications to the Keller--Segel Equation
\item
 Hamilton--Jacobi Equations and Weak KAM Theory
\item
 Free Boundary Problems
\item
 The Semi-Geostrophic Equations
\item
 Optimal Transport on Non-Compact Manifolds
\item
 Transport in Sub-Riemannian Geometry
\item
 Impact on Mathematics and Science
\item
 Awards and Recognition
\item
 Selected Major Publications
\item
 Ongoing Work and Future Directions
\item
 Conclusion
\item
 References and Further Reading
\end{enumerate}

\begin{center}\rule{0.5\linewidth}{0.5pt}\end{center}

\section{1. Introduction}\label{ch2018alessiofigalli:introduction}

Alessio Figalli was awarded the Fields Medal in 2018 for his
groundbreaking contributions to the theory of optimal transport and its
applications across partial differential equations, metric geometry, and
probability. At 34 years old, he had already produced over 150
publications spanning a notable range of topics in analysis,
geometry, and the calculus of variations.

Figalli's work is distinguished by its combination of technical power
and conceptual clarity. He has resolved fundamental problems in the
regularity theory of optimal transport maps, proven sharp quantitative
versions of classical geometric inequalities, and applied optimal
transport techniques to problems ranging from the Keller--Segel equation
in biology to the semi-geostrophic equations in meteorology.

His career trajectory is equally impressive: a PhD completed in a single
year under the supervision of two Fields Medalists (Ambrosio and
Villani), a rapid rise through academic positions in France, the United
States, and Switzerland, and a reputation as one of the most prolific
and influential mathematicians of his generation.

This report provides a detailed account of Figalli's important
contributions, the mathematical context in which they arose, the
techniques he developed, and their lasting impact.

\begin{center}\rule{0.5\linewidth}{0.5pt}\end{center}

\section{2. Biographical Background}\label{ch2018alessiofigalli:biographical-background}

\subsection{2.1 Early Life and
Education}\label{ch2018alessiofigalli:early-life-and-education}

Alessio Figalli was born on April 2, 1984, in Rome, Italy. He attended
the Liceo Classico Francesco Vivona in Rome, where he showed early
aptitude for mathematics. In 2002, he entered the Scuola Normale
Superiore di Pisa, one of Italy's most prestigious institutions for
advanced education. He earned his bachelor's and master's degrees in
mathematics from the University of Pisa in 2004 and 2006,
respectively---completing his master's a year ahead of schedule.

\subsection{2.2 Doctoral Studies}\label{ch2018alessiofigalli:doctoral-studies}

Figalli pursued his PhD under the joint supervision of Luigi Ambrosio at
the Scuola Normale Superiore di Pisa and Cedric Villani at the Ecole
Normale Superieure de Lyon. He completed his doctorate in 2007 in a
notable single year, with a thesis titled ``Optimal transportation
and action-minimizing measures.'' The thesis already contained results
of lasting significance, including work on the regularity of optimal
transport maps and connections to Lagrangian dynamics.

\subsection{2.3 Academic Career}\label{ch2018alessiofigalli:academic-career}

Figalli's academic career advanced with notable speed:

\begin{itemize}
\tightlist
\item
 \textbf{2007:} Appointed Charge de recherche at the French National
 Centre for Scientific Research (CNRS).
\item
 \textbf{2008:} Professeur Hadamard at the Ecole Polytechnique in
 Paris.
\item
 \textbf{2009:} Moved to the University of Texas at Austin as an
 associate professor.
\item
 \textbf{2011:} Promoted to full professor at UT Austin.
\item
 \textbf{2013:} Appointed to the R. L. Moore Chair at UT Austin.
\item
 \textbf{2016:} Moved to ETH Zurich as a chaired professor.
\item
 \textbf{2019:} Appointed Director of the Forschungsinstitut fur
 Mathematik (FIM) at ETH Zurich.
\end{itemize}

Figalli has supervised numerous doctoral students, including Maria
Colombo (now a professor at EPFL) and Robin Neumayer (now at
Northwestern University).

\begin{center}\rule{0.5\linewidth}{0.5pt}\end{center}

\section{3. Optimal Transport Theory: An
Overview}\label{ch2018alessiofigalli:optimal-transport-theory-an-overview}

\subsection{3.1 The Monge Problem}\label{ch2018alessiofigalli:the-monge-problem}

The optimal transport problem was first formulated by the French
mathematician Gaspard Monge in 1781. Given two probability distributions
--- a source distribution \(\mu\) and a target distribution \(\nu\) ---
the problem asks for a map \(T\) that transports \(\mu\) to \(\nu\)
while minimizing a cost functional:

\[
\inf_{T_\#\mu = \nu} \int c(x, T(x)) \, d\mu(x)
\]

When the cost is the squared Euclidean distance \(c(x,y) = |x-y|^2\),
the problem has a rich geometric structure. Transporting rubble to build
a fortification (Monge's original motivation) corresponds to a linear
cost in distance, but the squared distance cost is more natural for many
applications and has better mathematical properties.

\subsection{3.2 The Kantorovich
Formulation}\label{ch2018alessiofigalli:the-kantorovich-formulation}

In the 1940s, Leonid Kantorovich reformulated the problem in terms of
transport plans (joint probability distributions) rather than transport
maps. This relaxed formulation is a linear programming problem and
always has a solution. The Kantorovich formulation also introduces the
dual problem, which is fundamental for the analysis and computation of
optimal transport.

\subsection{3.3 The Brenier Theorem}\label{ch2018alessiofigalli:the-brenier-theorem}

A landmark result by Yann Brenier (1991) shows that for the quadratic
cost, the optimal transport map \(T\) is the gradient of a convex
function: \(T = \nabla \phi\). Moreover, \(\phi\) satisfies the
\textbf{Monge--Ampere equation}:

\[
\det(D^2 \phi(x)) = \frac{d\mu}{d\nu(\nabla \phi(x))}
\]

This connection between optimal transport and the Monge--Ampere equation
is the foundation of much of Figalli's work.

\subsection{3.4 Applications of Optimal
Transport}\label{ch2018alessiofigalli:applications-of-optimal-transport}

Optimal transport has applications in:

\begin{itemize}
\tightlist
\item
 \textbf{Partial differential equations:} The Wasserstein gradient flow
 formulation of evolution equations.
\item
 \textbf{Geometry:} The study of curvature and concentration of
 measure.
\item
 \textbf{Probability:} The theory of Wasserstein distances and
 coupling.
\item
 \textbf{Economics:} Matching problems and resource allocation.
\item
 \textbf{Machine learning:} Generative models and domain adaptation.
\item
 \textbf{Meteorology:} Semi-geostrophic equations for atmospheric
 flows.
\end{itemize}

\begin{center}\rule{0.5\linewidth}{0.5pt}\end{center}

\section{4. Regularity Theory for Optimal Transport
Maps}\label{ch2018alessiofigalli:regularity-theory-for-optimal-transport-maps}

\subsection{4.1 The Problem}\label{ch2018alessiofigalli:the-problem}

Given two probability densities \(f\) and \(g\) on \(\mathbb{R}^d\), the
optimal transport map \(T = \nabla \phi\) for the quadratic cost solves
the Monge--Ampere equation:

\[
\det(D^2 \phi(x)) = \frac{f(x)}{g(\nabla \phi(x))}
\]

A fundamental question is: \textbf{How regular is the optimal transport
map?} Under what conditions on \(f\) and \(g\) is \(\phi\) smooth
(\(C^\infty\)) or at least \(C^{1,\alpha}\)?

The regularity theory for optimal transport maps was pioneered by
Caffarelli (1990s), who established that if the domains are convex and
the densities are bounded and bounded away from zero, then \(T\) is
smooth in the interior. However, boundary regularity and the case of
non-convex domains remained challenging.

\subsection{4.2 Figalli's Contributions}\label{ch2018alessiofigalli:figallis-contributions}

Figalli made several fundamental contributions to the regularity theory:

\textbf{Global regularity:} With Guido De Philippis, Figalli proved an
optimal higher integrability result for the second derivatives of
solutions to the Monge--Ampere equation. Their theorem states that
\(\det D^2 \phi \in L^p_{\text{loc}}\) implies
\(D^2 \phi \in L^{p+\varepsilon}_{\text{loc}}\) for some
\(\varepsilon > 0\), improving previous bounds.

\textbf{Partial regularity:} In a landmark 2014 paper with De Philippis,
Figalli proved a partial regularity result for optimal transport maps
when the target density is only bounded. They showed that the singular
set of the optimal transport map has Hausdorff codimension at least 1
--- a dramatic improvement over previous results that only gave upper
bounds on the size of the singular set.

\textbf{Boundary regularity:} Figalli developed new techniques to handle
boundary regularity for optimal transport maps, particularly when the
domains are not convex. This work has applications to problems in
economics and urban planning.

\subsection{4.3 The Key Insight}\label{ch2018alessiofigalli:the-key-insight}

The main difficulty in optimal transport regularity is that the
Monge--Ampere equation is fully nonlinear and degenerate. Figalli's
approach combines:

\begin{itemize}
\tightlist
\item
 \textbf{Alexandrov estimates:} A priori bounds on the convex
 potential.
\item
 \textbf{Caffarelli's perturbation theory:} Relating the solution to
 the homogeneous Monge--Ampere equation.
\item
 \textbf{New estimates:} For the integrability of second derivatives,
 using delicate harmonic analysis and geometric measure theory.
\end{itemize}

\begin{center}\rule{0.5\linewidth}{0.5pt}\end{center}

\section{5. The Monge--Ampere Equation}\label{ch2018alessiofigalli:the-mongeampere-equation}

\subsection{5.1 Background}\label{ch2018alessiofigalli:background}

The Monge--Ampere equation is the prototypical fully nonlinear elliptic
equation:

\[
\det(D^2 u) = f(x, u, \nabla u)
\]

It appears in: - Optimal transport (as described above) - Differential
geometry (prescribed Gaussian curvature) - Fluid dynamics (the
semi-geostrophic equations) - Economics (utility maximization)

The regularity theory for Monge--Ampere equations is highly developed
but many fundamental questions remain open.

\subsection{5.2 Figalli's Contributions}\label{ch2018alessiofigalli:figallis-contributions-1}

Figalli and De Philippis proved several landmark results on the
structure of solutions to the Monge--Ampere equation:

\textbf{Higher integrability:} For a convex solution \(u\) of
\(\det D^2 u = f\) with \(f \in L^p\), they proved that
\(D^2 u \in L^{p+\varepsilon}\) for some \(\varepsilon > 0\). This
result is optimal and has important consequences for the regularity of
optimal transport maps.

\textbf{The Sobolev regularity of convex functions:} In a highly
influential paper, Figalli and De Philippis characterized the Sobolev
regularity of convex functions, proving that a convex function belongs
to \(W^{2,p}\) if and only if its Monge--Ampere measure is in \(L^p\).
This elegant characterization resolved a long-standing question.

\textbf{Partial regularity for Monge--Ampere type equations:} For a
class of equations that includes the optimal transport equation, Figalli
and De Philippis proved that the singular set of the solution has
codimension at least 1. This required new techniques because the
standard methods for fully nonlinear equations do not apply when the
right-hand side is only bounded.

\subsection{5.3 Convex Integration and
Non-Uniqueness}\label{ch2018alessiofigalli:convex-integration-and-non-uniqueness}

In joint work with De Philippis and others, Figalli has studied the
limits of the regularity theory. They showed that for certain
Monge--Ampere equations with rough right-hand sides, solutions can be
highly irregular, demonstrating that the regularity assumptions in the
classical theory are sharp.

\begin{center}\rule{0.5\linewidth}{0.5pt}\end{center}

\section{6. Partial Regularity for Monge--Ampere Type
Equations}\label{ch2018alessiofigalli:partial-regularity-for-mongeampere-type-equations}

\subsection{6.1 The Problem}\label{ch2018alessiofigalli:the-problem-1}

For a general fully nonlinear elliptic equation \(F(D^2 u) = f\), the
Evans--Krylov theory gives interior \(C^{2,\alpha}\) estimates when
\(f\) is smooth. When \(f\) is only bounded, the best one can hope for
is that the solution is \(C^{1,\alpha}\) outside a small singular set.

Proving that the singular set is small --- i.e., that \textbf{partial
regularity} holds --- is a major challenge. For optimal transport, the
equation is \(\det D^2 \phi = f\) with \(f\) possibly only bounded, and
the standard Evans--Krylov theory does not apply.

\subsection{6.2 Figalli and De Philippis's
Theorem}\label{ch2018alessiofigalli:figalli-and-de-philippiss-theorem}

In their breakthrough 2014 paper, Figalli and De Philippis proved:

\begin{quote}
\textbf{Theorem:} Let \(\Omega \subset \mathbb{R}^d\) be a bounded open
set, and let \(u: \Omega \to \mathbb{R}\) be a convex solution of
\(\det D^2 u = f\) with \(0 < \lambda \le f \le \Lambda\). Then there
exists a closed set \(\Sigma \subset \Omega\) such that: 1.
\(u \in C^{1,\alpha}_{\text{loc}}(\Omega \setminus \Sigma)\) for some
\(\alpha > 0\). 2. \(\dim_{\mathcal{H}}(\Sigma) \le d-1\) (the Hausdorff
dimension of \(\Sigma\) is at most \(d-1\)).
\end{quote}

In particular, the optimal transport map \(T = \nabla u\) is continuous
outside a set of dimension at most \(d-1\).

\subsection{6.3 Consequences}\label{ch2018alessiofigalli:consequences}

This result has important consequences:

\begin{itemize}
\tightlist
\item
 \textbf{Uniqueness of optimal maps:} The optimal transport map is
 uniquely determined and continuous almost everywhere.
\item
 \textbf{Stability:} Small perturbations of the densities lead to small
 perturbations of the transport map in \(C^1\) outside a small set.
\item
 \textbf{Numerical analysis:} The partial regularity result justifies
 the use of finite element methods for optimal transport.
\end{itemize}

\begin{center}\rule{0.5\linewidth}{0.5pt}\end{center}

\section{7. Quantitative Stability of Geometric and Functional
Inequalities}\label{ch2018alessiofigalli:quantitative-stability-of-geometric-and-functional-inequalities}

\subsection{7.1 Why Stability Matters}\label{ch2018alessiofigalli:why-stability-matters}

Many important inequalities in analysis and geometry are sharp only for
specific configurations. For example, the isoperimetric inequality
states that among all sets of given volume, the sphere minimizes
perimeter. The equality case is the sphere.

A \textbf{stability} (or \textbf{quantitative}) version of such an
inequality asks: if a set is almost optimal --- its perimeter is close
to minimizing --- must it be close to a sphere in some metric?
Understanding stability is essential for:

\begin{itemize}
\tightlist
\item
 Proving convergence of numerical methods.
\item
 Understanding the behavior of minimizing sequences.
\item
 Applications to evolution equations and geometric flows.
\end{itemize}

\subsection{7.2 Figalli's Approach}\label{ch2018alessiofigalli:figallis-approach}

Figalli developed powerful new techniques for proving quantitative
stability results for a wide range of inequalities. His approach
typically involves:

\begin{enumerate}
\def\labelenumi{\arabic{enumi}.}
\tightlist
\item
 \textbf{A compactness argument:} Show that a sequence of
 almost-optimal configurations converges to an optimal one.
\item
 \textbf{A linearization:} Analyze the second variation around the
 optimal configuration.
\item
 \textbf{Spectral analysis:} Use the eigenvalue of the relevant linear
 operator to obtain quantitative rates.
\end{enumerate}

This methodology was developed systematically in his work with Francesco
Maggi and Aldo Pratelli, and later extended by Figalli and many others.

\begin{center}\rule{0.5\linewidth}{0.5pt}\end{center}

\section{8. The Anisotropic Isoperimetric
Inequality}\label{ch2018alessiofigalli:the-anisotropic-isoperimetric-inequality}

\subsection{8.1 The Inequality}\label{ch2018alessiofigalli:the-inequality}

The \textbf{anisotropic isoperimetric inequality} generalizes the
classical isoperimetric inequality to account for direction-dependent
surface energies. Given a convex, positively homogeneous function
\(\Phi\) that defines a norm, the anisotropic perimeter of a set \(E\)
is:

\[
P_\Phi(E) = \int_{\partial E} \Phi(\nu_E) \, d\mathcal{H}^{d-1}
\]

where \(\nu_E\) is the outward normal. The inequality states that for
sets of given volume, the Wulff shape (the unit ball of the dual norm)
minimizes the anisotropic perimeter.

\subsection{8.2 Quantitative Stability}\label{ch2018alessiofigalli:quantitative-stability}

In a well-known 2010 paper with Francesco Maggi and Aldo Pratelli,
Figalli proved a sharp quantitative version of the anisotropic
isoperimetric inequality:

\begin{quote}
\textbf{Theorem:} There exists a constant \(C = C(d, \Phi) > 0\) such
that for any set \(E \subset \mathbb{R}^d\) of finite perimeter, \[
P_\Phi(E) - P_\Phi(W) \ge C \, \alpha(E)^2
\] where \(W\) is the Wulff shape of the same volume as \(E\), and
\(\alpha(E)\) is the \(L^1\)-distance (or any natural distance) between
\(E\) and \(W\).
\end{quote}

The exponent 2 is optimal, meaning the inequality is as sharp as
possible.

\subsection{8.3 Significance}\label{ch2018alessiofigalli:significance}

This result was a major advance in the field of geometric measure
theory. Prior results gave stability with exponent 1 (a linear rate),
which is not optimal. Figalli, Maggi, and Pratelli's proof combined:

\begin{itemize}
\tightlist
\item
 The selection principle: A compactness argument to extract a
 convergent subsequence.
\item
 A new inequality relating the perimeter deficit to the asymmetry
 index.
\item
 Optimal transport techniques to connect the two.
\end{itemize}

The anisotropic isoperimetric inequality and its stability have
applications to:

\begin{itemize}
\tightlist
\item
 \textbf{Crystallization:} The Wulff shape describes the equilibrium
 crystal shape.
\item
 \textbf{Surface tension:} The direction-dependent surface energy
 appears in materials science.
\item
 \textbf{Phase transitions:} The stability estimate controls the rate
 of coarsening.
\end{itemize}

\begin{center}\rule{0.5\linewidth}{0.5pt}\end{center}

\section{9. Stability of the Gagliardo--Nirenberg and Log-Sobolev
Inequalities}\label{ch2018alessiofigalli:stability-of-the-gagliardonirenberg-and-log-sobolev-inequalities}

\subsection{9.1 Functional Inequalities}\label{ch2018alessiofigalli:functional-inequalities}

The Gagliardo--Nirenberg inequalities and the logarithmic Sobolev
inequality are fundamental tools in analysis and probability. They
control the \(L^p\) norm of a function in terms of its gradient and its
\(L^q\) norm. The logarithmic Sobolev inequality, in particular, is
central to the theory of Markov processes and statistical mechanics.

\subsection{9.2 Quantitative Stability with
Carlen}\label{ch2018alessiofigalli:quantitative-stability-with-carlen}

In a 2013 paper with Eric Carlen, Figalli proved sharp quantitative
stability for a family of Gagliardo--Nirenberg inequalities and the
logarithmic Sobolev inequality. Their main theorem states:

\begin{quote}
If a function nearly achieves equality in one of these inequalities,
then it must be close to an extremal function (a Gaussian or a幂函数),
with a quantitative rate that is quadratic in the distance.
\end{quote}

\subsection{9.3 Application to the Keller--Segel
Equation}\label{ch2018alessiofigalli:application-to-the-kellersegel-equation}

The stability result for the logarithmic Sobolev inequality has an
important application: it provides a \textbf{quantitative rate of
convergence} for the critical mass Keller--Segel equation.

The Keller--Segel equation models chemotaxis --- the movement of
organisms (like bacteria) in response to chemical signals. In the
critical mass case, solutions are global and converge to a steady state.
Figalli and Carlen's stability estimate gives an explicit rate of
convergence, solving a problem that had been open for over a decade.

\begin{center}\rule{0.5\linewidth}{0.5pt}\end{center}

\section{10. Applications to the Keller--Segel
Equation}\label{ch2018alessiofigalli:applications-to-the-kellersegel-equation}

\subsection{10.1 The Equation}\label{ch2018alessiofigalli:the-equation}

The (Patlak--)Keller--Segel system describes the evolution of a cell
density \(\rho(x,t)\) and a chemical concentration \(c(x,t)\):

\[
\begin{cases}
\partial_t \rho = \Delta \rho - \nabla \cdot (\rho \nabla c) \\
0 = \Delta c + \rho
\end{cases}
\]

The first equation is a conservation law with diffusion and chemotactic
drift; the second is the Poisson equation for the chemical
concentration.

The Keller--Segel equation exhibits a critical mass phenomenon: -
\textbf{Subcritical mass:} Solutions exist globally and become extinct.
- \textbf{Critical mass:} Solutions exist globally and converge to a
steady state. - \textbf{Supercritical mass:} Solutions blow up in finite
time.

\subsection{10.2 Figalli's Contribution}\label{ch2018alessiofigalli:figallis-contribution}

Figalli, together with Eric Carlen, used optimal transport techniques to
study the critical mass case. Their main result:

\begin{itemize}
\tightlist
\item
 \textbf{Quantitative convergence:} The solution converges to the
 stationary state at an explicit algebraic rate. The rate is determined
 by the spectral gap of a certain linear operator.
\item
 \textbf{Sharp constants:} The rate is optimal, in the sense that it
 cannot be improved.
\end{itemize}

The proof uses the \textbf{Wasserstein gradient flow} structure of the
Keller--Segel equation. The equation can be written as:

\[
\partial_t \rho = \nabla \cdot (\rho \nabla (E'(\rho) - c))
\]

which is a gradient flow for the free energy

\[
F(\rho) = \int \rho \log \rho \, dx - \frac{1}{2} \iint \rho(x) G(x-y) \rho(y) \, dx \, dy
\]

where \(G\) is the Newtonian potential. The logarithmic Sobolev
inequality controls the dissipation of this free energy.

\begin{center}\rule{0.5\linewidth}{0.5pt}\end{center}

\section{11. Hamilton--Jacobi Equations and Weak KAM
Theory}\label{ch2018alessiofigalli:hamiltonjacobi-equations-and-weak-kam-theory}

\subsection{11.1 Background}\label{ch2018alessiofigalli:background-1}

Hamilton--Jacobi equations are first-order fully nonlinear PDEs of the
form:

\[
\partial_t u + H(x, \nabla u) = 0
\]

They arise in classical mechanics, geometric optics, and optimal control
theory. \textbf{Weak KAM theory}, developed by Albert Fathi and others,
studies the long-time behavior of solutions to Hamilton--Jacobi
equations and their connection to the dynamics of the Hamiltonian
system.

\subsection{11.2 Figalli's
Contributions}\label{ch2018alessiofigalli:figallis-contributions-2}

Figalli has made several contributions to this area:

\textbf{Generic hyperbolicity of Aubry sets:} In joint work with Gonzalo
Contreras and Ludovic Rifford, Figalli proved that for a generic Tonelli
Hamiltonian on a compact surface, the Aubry set --- the set of ``most
recurrent'' orbits --- is a hyperbolic set. This resolves a
long-standing conjecture in dynamical systems.

\textbf{Optimal transport and weak KAM:} Figalli's PhD thesis
established deep connections between optimal transport theory and weak
KAM theory. The key idea is that the optimal transport problem for a
Lagrangian cost is equivalent to finding solutions of the
Hamilton--Jacobi equation. Specifically, the value function of the
optimal transport problem:

\[
u(x) = \inf_{y} [c(x,y) - v(y)]
\]

satisfies the Hamilton--Jacobi equation in a weak sense.

\textbf{Action-minimizing measures:} Figalli studied the structure of
action-minimizing measures for Lagrangian systems, proving their
regularity and establishing connections with Mather theory.

\subsection{11.3 The Weak KAM/OT
Dictionary}\label{ch2018alessiofigalli:the-weak-kamot-dictionary}

The connection between optimal transport and weak KAM theory has become
an important tool in both fields. The correspondence is:

{\def\LTcaptype{none} % do not increment counter
\begin{longtable}[]{@{}ll@{}}
\toprule\noalign{}
Optimal Transport & Weak KAM Theory \\
\midrule\noalign{}
\endhead
\bottomrule\noalign{}
\endlastfoot
Cost function \(c(x,y)\) & Lagrangian \(L(x,v)\) \\
Kantorovich potential & Weak KAM solution \\
Optimal transport map & Legendre transform of the Lagrangian flow \\
Wasserstein distance & Mather's action functional \\
\end{longtable}
}

\begin{center}\rule{0.5\linewidth}{0.5pt}\end{center}

\section{12. Free Boundary Problems}\label{ch2018alessiofigalli:free-boundary-problems}

\subsection{12.1 What Are Free Boundary
Problems?}\label{ch2018alessiofigalli:what-are-free-boundary-problems}

A free boundary problem is a PDE where the domain of the equation is
unknown and must be determined as part of the solution. The ``free
boundary'' is the boundary between two regions where different equations
hold. Examples include: - The Stefan problem (melting ice) - The
obstacle problem (a membrane pressed against an obstacle) - Optimal
stopping in finance

\subsection{12.2 Figalli's
Contributions}\label{ch2018alessiofigalli:figallis-contributions-3}

Figalli has worked on several free boundary problems, particularly those
arising from optimal transport:

\textbf{The optimal partial transport problem:} This is a variant of
optimal transport where only part of the mass is transported. The free
boundary separates the region where mass is moved from the region where
it stays put. Figalli proved regularity results for the free boundary in
this problem.

\textbf{The obstacle problem for Monge--Ampere:} In joint work with De
Philippis and others, Figalli studied the obstacle problem for the
Monge--Ampere equation, proving regularity of the free boundary.

\textbf{Applications to economics:} The optimal partial transport
problem has applications to matching problems with capacity constraints,
where the free boundary represents the boundary of the ``active''
market.

\begin{center}\rule{0.5\linewidth}{0.5pt}\end{center}

\section{13. The Semi-Geostrophic
Equations}\label{ch2018alessiofigalli:the-semi-geostrophic-equations}

\subsection{13.1 Background}\label{ch2018alessiofigalli:background-2}

The semi-geostrophic equations are a system of PDEs used in meteorology
and oceanography to model large-scale atmospheric and oceanic flows.
They approximate the full Euler equations under the assumption of strong
rotation (the Coriolis force dominates).

The equations can be written in terms of a potential \(\phi\) that
satisfies a Monge--Ampere type equation. The connection between the
semi-geostrophic equations and optimal transport was discovered by
Cullen and Purser in the 1980s.

\subsection{13.2 Figalli's Contribution}\label{ch2018alessiofigalli:figallis-contribution-1}

With Guido De Philippis, Figalli made a breakthrough in understanding
the semi-geostrophic equations by proving:

\begin{itemize}
\tightlist
\item
 \textbf{Well-posedness:} Existence and uniqueness of solutions for a
 class of initial data.
\item
 \textbf{Regularity:} Solutions are smooth for smooth initial data.
\item
 \textbf{The Monge--Ampere connection:} The semi-geostrophic potential
 satisfies a Monge--Ampere equation that can be interpreted as the
 optimal transport between the initial density and the evolving
 density.
\end{itemize}

The main result is that the semi-geostrophic equations can be solved
using optimal transport: at each time, the geostrophic potential is the
solution of an optimal transport problem between the current density and
the Lagrangian reference density.

This connection provides an effective tool for both the theoretical
analysis and the numerical simulation of atmospheric flows.

\begin{center}\rule{0.5\linewidth}{0.5pt}\end{center}

\section{14. Optimal Transport on Non-Compact
Manifolds}\label{ch2018alessiofigalli:optimal-transport-on-non-compact-manifolds}

\subsection{14.1 The Problem}\label{ch2018alessiofigalli:the-problem-2}

The classical theory of optimal transport was developed primarily on
compact manifolds or on \(\mathbb{R}^d\) with certain growth conditions.
Extending the theory to non-compact manifolds requires additional care,
as the cost function may not be finite for all pairs of points.

\subsection{14.2 Figalli and Fathi's
Theorem}\label{ch2018alessiofigalli:figalli-and-fathis-theorem}

In a 2010 paper with Albert Fathi, Figalli extended the optimal
transport theory to general complete Riemannian manifolds. They proved:

\begin{quote}
\textbf{Theorem:} On any complete Riemannian manifold, for costs of the
form \(c(x,y) = d(x,y)^r\) with \(r > 1\), the optimal transport problem
is well-posed. The optimal transport map exists and is unique, and it is
given by the gradient of a \(c\)-convex function.
\end{quote}

This result holds without any curvature restrictions, and it includes
the important case of the quadratic cost \(c(x,y) = d(x,y)^2\).

\subsection{14.3 Significance}\label{ch2018alessiofigalli:significance-1}

Before this work, the existence of optimal transport maps on non-compact
manifolds was known only under restrictive assumptions (bounded
geometry, non-negative curvature, etc.). Figalli and Fathi's result
showed that optimal transport is a robust theory that works in full
generality.

\begin{center}\rule{0.5\linewidth}{0.5pt}\end{center}

\section{15. Transport in Sub-Riemannian
Geometry}\label{ch2018alessiofigalli:transport-in-sub-riemannian-geometry}

\subsection{15.1 Sub-Riemannian Geometry}\label{ch2018alessiofigalli:sub-riemannian-geometry}

Sub-Riemannian geometry generalizes Riemannian geometry by allowing only
horizontal directions. The Heisenberg group is the simplest example. In
sub-Riemannian geometry, the distance between points can be very
different from the Euclidean distance, and the geometry is more
challenging.

\subsection{15.2 Figalli's
Contributions}\label{ch2018alessiofigalli:figallis-contributions-4}

With Ludovic Rifford, Figalli studied optimal transport on
sub-Riemannian manifolds:

\begin{itemize}
\tightlist
\item
 \textbf{Existence and uniqueness:} They proved that for the Heisenberg
 group and more general Carnot groups, the optimal transport problem
 for the squared sub-Riemannian distance is well-posed.
\item
 \textbf{Regularity:} Unlike the Riemannian case, the optimal transport
 map may be discontinuous in sub-Riemannian geometry. Figalli and
 Rifford characterized the singular set.
\item
 \textbf{Absolute continuity of Wasserstein geodesics:} With Nicolas
 Juillet, Figalli proved that Wasserstein geodesics in the Heisenberg
 group are absolutely continuous, a non-trivial result due to the
 degenerate geometry.
\end{itemize}

This work opened up the study of optimal transport in sub-Riemannian
geometry, which is now an active research area with applications to
control theory and robotics.

\begin{center}\rule{0.5\linewidth}{0.5pt}\end{center}

\section{16. Impact on Mathematics and
Science}\label{ch2018alessiofigalli:impact-on-mathematics-and-science}

\subsection{16.1 Impact on Optimal Transport
Theory}\label{ch2018alessiofigalli:impact-on-optimal-transport-theory}

Figalli is a leading figure in optimal
transport theory. His work on regularity, partial regularity, and
stability has substantially advanced the field. The Figalli--De
Philippis partial regularity theorem is now a standard result taught in
advanced courses on optimal transport.

\subsection{16.2 Impact on PDE Theory}\label{ch2018alessiofigalli:impact-on-pde-theory}

Figalli's results on the Monge--Ampere equation have influenced the
broader theory of fully nonlinear elliptic PDEs. His characterization of
Sobolev regularity for convex functions is a foundational result, and
his work on the semi-geostrophic equations has opened new connections
between optimal transport and fluid dynamics.

\subsection{16.3 Impact on Geometric Measure
Theory}\label{ch2018alessiofigalli:impact-on-geometric-measure-theory}

The quantitative stability results for isoperimetric inequalities,
proved with Maggi and Pratelli, have become a cornerstone of geometric
measure theory. They have been extended to many other settings
(Riemannian manifolds, metric measure spaces) and have applications to:

\begin{itemize}
\tightlist
\item
 \textbf{Phase transitions:} The stability estimate controls the rate
 of convergence in Allen--Cahn dynamics.
\item
 \textbf{Minimal surfaces:} The quantitative estimate is used in the
 study of almost-minimizers.
\item
 \textbf{Crystallization:} The anisotropic version describes the
 approach to equilibrium in crystal growth.
\end{itemize}

\subsection{16.4 Impact on Probability and
Statistics}\label{ch2018alessiofigalli:impact-on-probability-and-statistics}

Figalli's work on optimal transport in probability and statistics has
influenced:

\begin{itemize}
\tightlist
\item
 \textbf{Wasserstein gradient flows:} Used in the analysis of sampling
 algorithms and Bayesian inference.
\item
 \textbf{Transportation of measure:} Applied in machine learning for
 generative models and domain adaptation.
\item
 \textbf{Stability of functional inequalities:} Used in the study of
 convergence rates for Markov chains.
\end{itemize}

\subsection{16.5 Impact on Fluid Dynamics and
Meteorology}\label{ch2018alessiofigalli:impact-on-fluid-dynamics-and-meteorology}

The connection between optimal transport and the semi-geostrophic
equations, developed by Figalli and De Philippis, has provided new tools
for the analysis of atmospheric flows. The optimal transport
interpretation allows:

\begin{itemize}
\tightlist
\item
 \textbf{Global existence:} Results for the semi-geostrophic equations
 with rough initial data.
\item
 \textbf{Numerical methods:} Structure-preserving discretizations based
 on optimal transport.
\item
 \textbf{Modeling:} A deeper understanding of front formation and
 weather patterns.
\end{itemize}

\subsection{16.6 Mentorship and
Leadership}\label{ch2018alessiofigalli:mentorship-and-leadership}

Figalli has supervised numerous PhD students and postdocs who have gone
on to become leading researchers. His students include Maria Colombo
(Fields Medalist of the future?), Robin Neumayer, and many others.

Since 2019, he has been Director of the FIM (Forschungsinstitut fur
Mathematik) at ETH Zurich, one of Europe's leading mathematics research
institutes. Under his leadership, FIM has continued to attract top
researchers and organize influential programs.

\begin{center}\rule{0.5\linewidth}{0.5pt}\end{center}

\section{17. Awards and Recognition}\label{ch2018alessiofigalli:awards-and-recognition}

\textbf{2024:} - Golden Owl for exceptional teaching, ETH Zurich - UIMP
Medal

\textbf{2023:} - Frontiers of Science Award

\textbf{2020:} - Falling Walls Award (Science and Innovation Management)

\textbf{2019:} - Doctorate Honoris Causa, Polytechnic University of
Catalonia

\textbf{2018:} - Fields Medal (International Mathematical Union) -
Doctorate Honoris Causa, Universite Cote d'Azur

\textbf{2017:} - Feltrinelli Prize for Mathematics (Accademia dei
Lincei)

\textbf{2015:} - Stampacchia Gold Medal

\textbf{2014:} - Invited Speaker, International Congress of
Mathematicians, Seoul

\textbf{2012:} - EMS Prize (European Mathematical Society) -
Peccot-Vimont Prize and Cours Peccot, College de France

\textbf{2011:} - Peccot Lectures, College de France

\begin{center}\rule{0.5\linewidth}{0.5pt}\end{center}

\section{18. Selected Major
Publications}\label{ch2018alessiofigalli:selected-major-publications}

\begin{enumerate}
\def\labelenumi{\arabic{enumi}.}
\item
 \textbf{``A quantitative isoperimetric inequality in
 \(\mathbb{R}^d\)''} (with F. Maggi and A. Pratelli). \emph{Journal of
 the American Mathematical Society} 23(1), 171--202, 2010. ---
 \emph{Sharp quantitative stability for the isoperimetric inequality,
 with optimal exponent 2.}
\item
 \textbf{``Partial regularity of optimal transport maps''} (with G. De
 Philippis). \emph{Annals of Mathematics} 180(1), 193--225, 2014. ---
 \emph{Proof that the singular set of an optimal transport map has
 codimension at least 1.}
\item
 \textbf{``\(W^{2,1}\) regularity for solutions of the Monge--Ampere
 equation''} (with G. De Philippis). \emph{Inventiones Mathematicae}
 192(1), 55--69, 2013. --- \emph{Characterization of Sobolev regularity
 for convex solutions of the Monge--Ampere equation.}
\item
 \textbf{``Stability for the logarithmic Sobolev inequality and the
 Keller--Segel system''} (with E. Carlen). \emph{Journal of the
 American Mathematical Society} 28(4), 1073--1108, 2015. ---
 \emph{Quantitative stability for functional inequalities and
 application to chemotaxis.}
\item
 \textbf{``Optimal transportation on non-compact manifolds''} (with A.
 Fathi). \emph{Israel Journal of Mathematics} 175(1), 1--59, 2010. ---
 \emph{Existence and uniqueness of optimal transport maps on complete
 Riemannian manifolds.}
\item
 \textbf{``Mass Transportation on Sub-Riemannian Manifolds''} (with L.
 Rifford). \emph{Geometric and Functional Analysis} 20(1), 124--159,
 2010. --- \emph{Optimal transport theory for the Heisenberg group and
 Carnot groups.}
\item
 \textbf{``The optimal partial transport problem.''} \emph{Archive for
 Rational Mechanics and Analysis} 195(2), 533--560, 2010. ---
 \emph{Free boundary problem for optimal transport with partial mass.}
\item
 \textbf{``A note on the regularity of the free boundaries in the
 optimal partial transport problem.''} \emph{Rendiconti del Circolo
 Matematico di Palermo} 58(2), 283--286, 2009.
\item
 \textbf{``Generic hyperbolicity of Aubry sets on surfaces''} (with G.
 Contreras and L. Rifford). \emph{Inventiones Mathematicae} 200(1),
 201--257, 2015. --- \emph{Proof of generic hyperbolicity in weak KAM
 theory.}
\item
 \textbf{``The Monge--Ampere equation and its link to optimal
 transportation.''} \emph{Bulletin of the American Mathematical
 Society} 55(4), 391--419, 2018. --- \emph{Survey article connecting
 Monge--Ampere theory with optimal transport.}
\item
 \textbf{``Regularity of optimal transport maps and partial
 differential equations.''} \emph{Proceedings of the International
 Congress of Mathematicians} 2014. --- \emph{ICM invited lecture
 surveying Figalli's work.}
\item
 \textbf{``An introduction to optimal transport and Wasserstein
 gradient flows.''} \emph{Springer Lecture Notes}, 2024. ---
 \emph{Comprehensive lecture notes on modern optimal transport.}
\end{enumerate}

\begin{center}\rule{0.5\linewidth}{0.5pt}\end{center}

\section{19. Ongoing Work and Future
Directions}\label{ch2018alessiofigalli:ongoing-work-and-future-directions}

\subsection{19.1 Regularity Theory}\label{ch2018alessiofigalli:regularity-theory}

Figalli continues to work on the regularity theory for optimal transport
and Monge--Ampere equations. Current directions include:

\begin{itemize}
\tightlist
\item
 \textbf{Boundary regularity:} Complete understanding of the behavior
 of optimal transport maps at the boundary of the domain.
\item
 \textbf{Higher codimension singularities:} Understanding when the
 singular set has dimension exactly \(d-1\).
\item
 \textbf{Non-quadratic costs:} Extending the regularity theory to costs
 beyond the quadratic case.
\end{itemize}

\subsection{19.2 Free Boundary Problems}\label{ch2018alessiofigalli:free-boundary-problems-1}

Figalli is working on several free boundary problems, including:

\begin{itemize}
\tightlist
\item
 \textbf{The Stefan problem:} Optimal transport methods for the
 classical melting problem.
\item
 \textbf{The Obstacle problem:} Higher regularity for the free
 boundary.
\item
 \textbf{Two-phase problems:} Free boundary problems where both phases
 are active.
\end{itemize}

\subsection{19.3 Machine Learning and Optimal
Transport}\label{ch2018alessiofigalli:machine-learning-and-optimal-transport}

Figalli has recently become interested in the connections between
optimal transport and machine learning:

\begin{itemize}
\tightlist
\item
 \textbf{Generative models:} Wasserstein GANs and diffusion models.
\item
 \textbf{Domain adaptation:} Transporting data distributions between
 domains.
\item
 \textbf{Uncertainty quantification:} Using Wasserstein distances to
 measure model uncertainty.
\end{itemize}

\subsection{19.4 Random Matrix Theory}\label{ch2018alessiofigalli:random-matrix-theory}

Figalli has recently begun working on random matrix theory, using
optimal transport techniques to study:

\begin{itemize}
\tightlist
\item
 \textbf{Eigenvalue distributions:} The convergence of empirical
 eigenvalue distributions to the semicircle law.
\item
 \textbf{Dyson's Brownian motion:} The dynamics of eigenvalues under
 random perturbations.
\item
 \textbf{Optimal transport and free probability:} Connections between
 optimal transport and Voiculescu's free probability theory.
\end{itemize}

\subsection{19.5 Teaching and Mentoring}\label{ch2018alessiofigalli:teaching-and-mentoring}

Figalli remains closely committed to teaching and mentoring. He received
the ETH Golden Owl for exceptional teaching in 2024, and he continues to
supervise PhD students and postdocs. He has given numerous advanced
lecture series, including the Nachdiplom Lectures at ETH Zurich.

\begin{center}\rule{0.5\linewidth}{0.5pt}\end{center}

\section{20. Conclusion}\label{ch2018alessiofigalli:conclusion}

Alessio Figalli's reception of the Fields Medal in 2018 recognized a
body of work that had already transformed several areas of mathematics.
His contributions to the regularity theory of optimal transport ---
including the partial regularity theorem with De Philippis and the
higher integrability results --- are fundamental advances. His
quantitative stability results for geometric and functional
inequalities, particularly the anisotropic isoperimetric inequality and
the logarithmic Sobolev inequality, are landmark results that have found
applications across analysis, geometry, and probability.

Figalli's work is characterized by its substantial breadth and depth.
He moves effortlessly between abstract theory and concrete applications,
from the deep geometric measure theory of isoperimetric problems to the
applied PDE analysis of the Keller--Segel and semi-geostrophic
equations. His technical power is matched by his conceptual clarity: his
papers are known for being well-written and accessible, and his lectures
are models of expository excellence.

At just 42 years old as of 2026, Figalli is at the height of his
mathematical powers. As Director of FIM at ETH Zurich, he continues to
shape the direction of research in analysis and PDEs. His ongoing work
on optimal transport, free boundary problems, machine learning, and
random matrix theory promises to yield further breakthroughs in the
years to come.

\begin{center}\rule{0.5\linewidth}{0.5pt}\end{center}

\section{21. References and Further
Reading}\label{ch2018alessiofigalli:references-and-further-reading}

\subsection{Official Sources}\label{ch2018alessiofigalli:official-sources}

\begin{itemize}
\tightlist
\item
 IMU Fields Medal 2018 Announcement:
 https://www.mathunion.org/imu-awards/fields-medal/fields-medals-2018
\item
 Fields Medal Citation for Alessio Figalli:
 https://www.mathunion.org/fileadmin/IMU/Prizes/Fields/2018/Figalli-Citation.pdf
\item
 Jackson, A. ``The Work of Alessio Figalli.'' ICM 2018 Laudatio.
\end{itemize}

\subsection{Biographical and Overview
Articles}\label{ch2018alessiofigalli:biographical-and-overview-articles}

\begin{itemize}
\tightlist
\item
 Klarreich, E. ``A Traveler Who Finds Stability in the Natural World.''
 \emph{Quanta Magazine}, August 1, 2018.
\item
 Figalli, A. ``The Monge--Ampere equation and its link to optimal
 transportation.'' \emph{Bulletin of the AMS} 55(4), 391--419, 2018.
\item
 Castelvecchi, D. ``Number-theory prodigy among winners of most coveted
 prize in mathematics.'' \emph{Nature}, August 1, 2018.
\end{itemize}

\subsection{Key Technical References}\label{ch2018alessiofigalli:key-technical-references}

\begin{itemize}
\tightlist
\item
 Figalli, A., Maggi, F., and Pratelli, A. ``A quantitative
 isoperimetric inequality in \(\mathbb{R}^d\).'' \emph{JAMS} 23 (2010),
 171--202.
\item
 De Philippis, G. and Figalli, A. ``Partial regularity of optimal
 transport maps.'' \emph{Annals of Mathematics} 180 (2014), 193--225.
\item
 De Philippis, G. and Figalli, A. ``\(W^{2,1}\) regularity for
 solutions of the Monge--Ampere equation.'' \emph{Inventiones
 Mathematicae} 192 (2013), 55--69.
\item
 Carlen, E. and Figalli, A. ``Stability for the logarithmic Sobolev
 inequality and the Keller--Segel system.'' \emph{JAMS} 28 (2015),
 1073--1108.
\item
 Fathi, A. and Figalli, A. ``Optimal transportation on non-compact
 manifolds.'' \emph{Israel Journal of Mathematics} 175 (2010), 1--59.
\item
 Figalli, A. and Rifford, L. ``Mass Transportation on Sub-Riemannian
 Manifolds.'' \emph{Geometric and Functional Analysis} 20 (2010),
 124--159.
\item
 Contreras, G., Figalli, A., and Rifford, L. ``Generic hyperbolicity of
 Aubry sets on surfaces.'' \emph{Inventiones Mathematicae} 200 (2015),
 201--257.
\item
 Villani, C. \emph{Optimal Transport: Old and New.} Springer, 2009.
\item
 Santambrogio, F. \emph{Optimal Transport for Applied Mathematicians.}
 Birkhauser, 2015.
\end{itemize}

\begin{center}\rule{0.5\linewidth}{0.5pt}\end{center}

\emph{(End of report)}
